\chapter{测试概述}\label{ch:overview}

\section{测试结论}

本文档以顶层模块 \code{AFDX_top} 及其全部 \RtlFileCount{} 个 RTL 源文件（\DesignUnitCount{} 个设计单元）为被测对象，验证赛题及官方 1022/1026/1030 答疑要求的 5 项核心功能、4 项关键技术指标、2B-VLID 顶层流封装、1--\MaxMessageBytes{}~B IPv4 分片与重组、代码覆盖率及综合与布局布线后门级网表行为。测试共编排 \RegCount{} 项 RTL 自动化回归测试平台、顶层 63-VL 满载规格测试套件、2 项配置隔离验证及 4 项 VQM 门级后仿真，覆盖 52 个测试点，各项功能与性能指标均达到或优于赛题规范，汇总结论见表~\ref{tab:summary}，测试层次与规模见表~\ref{tab:summary-scale}。

\begingroup\renewcommand{\tablearraystretch}{1.06}
\begin{longtable}{@{}>{\raggedright\arraybackslash}p{2.0cm}>{\raggedright\arraybackslash}p{2.5cm}>{\raggedright\arraybackslash}p{\dimexpr\linewidth-5.6cm-6\tabcolsep\relax}>{\centering\arraybackslash}p{1.1cm}@{}}
\caption{赛题要求与测试结论}\label{tab:summary}\\
\longtablehead{4}{\thead{赛题条目} & \thead{要求} & \thead{测试结果} & \thead{结论}}
\multicolumn{4}{@{}l}{\heiti 功能} \\
ARINC 664 通信 & 支持 P7 规定的通信功能 & UDP/IP 封装、VL 配置、BAG 调度、1026 顶层 2B-VLID 流包装及 1030 超长消息 IPv4 分片/重组双网交付全部正确（\ref{sec:f-comm}、\ref{sec:f-frag}~节） & 符合 \\
高完整性通信 & 必须支持 & 收发 14 类异常均准确识别、分类和计数；冗余窗口内双网副本严格单次交付，去重后按序重组（\ref{sec:f-ic}~节） & 符合 \\
SNMP 网络管理 & 支持 & v1/v2c 的 Get、GetNext、Set 与 v2 Trap 编解码正确，MIB 读回值与硬件寄存器一致（\ref{sec:f-snmp}~节） & 符合 \\
615A 软件数据加载 & 支持 & 单文件上传后片内 RAM 数据与介质定义下载（Media-Defined Download, MDD）回传文件均逐字节一致；差错与丢包按协议重传或终止（\ref{sec:f-615a}~节） & 符合 \\
RTC 时钟同步 & 支持 & 空载/满载同步误差分别为 \RtcIdle{}/\RtcLoad{}~ns；应答丢失后自动重发恢复，非法或过期应答不误触发时钟校准（\ref{sec:f-rtc}~节） & 符合 \\
\midrule
\multicolumn{4}{@{}l}{\heiti 技术指标} \\
传输速率 & 100~Mbps & 125~MHz 八位 GMII 下，顶层 \code{AFDX_top} 声明满载规格 63 条 VL 单向有效吞吐达 \ThruSixtyThree{}~Mbps（双向 \ThruSixtyThreeBi{}~Mbps）；超规格压力工况 77/81 条 VL 单向达 \Thru{}/\ThruEightyOne{}~Mbps，变长高负载单向达 \ThruVarHigh{}~Mbps（\ref{sec:p-rate}~节） & 符合 \\
帧格式 & 单帧载荷 $\le$1471~B、帧长 $\le$1518~B，支持 1--8192~B 分片及 1--255 序号 & 单帧 1--1471~B 载荷、64--1518~B 帧长逐字节核对正确；1472--\MaxMessageBytes{}~B 消息按 1472~B IPv4 分片并在去重后按序重组；VLID 映射、首帧 SN$=0$ 及 1--255 回绕正确（\ref{sec:p-frame}~节） & 符合 \\
时延与抖动 & 端到端 $<$1--2~ms，抖动不超过几十微秒 & 顶层 63 条 VL 稳态末字节最大时延 \LatSixtyThreePMax{}~µs（逐周期抖动 \JitSixtyThreePer{}~µs，两拍突发最大 \LatSixtyThreeBTwo{}~µs / 抖动 \JitSixtyThreeBTwo{}~µs）；单帧基线最大 \Lat{}~µs、同 VL 抖动 \JitSameVl{}~µs；81 条最长帧突发排空 \BurstEightOne{}~µs（\ref{sec:p-lat}~节） & 符合 \\
VL 支持能力 & 128 条发送、4096 条接收 & 128 条发送 VL 全量并发发送通过；4096 条接收 VL 全表配置并逐条完成双网去重交付（\ref{sec:p-vl}~节） & 符合 \\
\midrule
时序与网表 & 125~MHz 静态时序收敛 & 布局布线后（\code{PRES 1.30 + tiles_1030.edc}，\ClkUncert{}~ns 不确定度）净建立裕量 +\RoutedSetup~ns、保持裕量 +\RoutedHold~ns，\countsep{\RoutedEndpoints} 端点 0 违例；4 项 VQM 门级网表仿真与 RTL 一致（第~\ref{ch:netlist}~章） & 符合 \\
\end{longtable}\endgroup

\begin{table}[H]\centering\tablestyle
\caption{测试规模}\label{tab:summary-scale}
\begin{tabularx}{\linewidth}{@{}lLc@{}}
\toprule
\thead{测试} & \thead{规模} & \thead{结果} \\
\midrule
RTL 自动化回归 & \RegCount{} 项用例（\RegFastCount{} 项快速回归 + 2 项高负载/RTC + 2 项长时系统性能参考），自检 \countsep{\RegChecks} 次，独立核对 \countsep{\RegIndep} 项 & 全部通过 \\
63-VL 满载规格与配置隔离 & 顶层 63 条 VL 7 组定长基准工况 + 混合分片工况（1472/3000/8192~B，12,206 条消息）+ 2 项配置隔离测试 & 全部通过 \\
eLinx 行为仿真 & 双向多帧长双网通信测试 SYS-09（\code{tb_elinx}，\texttt{checks=\ElinxChecks{}, errors=0}） & 通过 \\
综合后门级仿真 & 4 项 VQM 门级网表测试（NET-01 至 NET-04：eLinx 双端互发、五功能逐字节比对、1030 分片窗口、1026 收发包装） & 全部通过 \\
\bottomrule
\end{tabularx}
\end{table}

\section{被测对象与测试依据}

被测对象为 \code{AFDX_top} 顶层及其全部 \RtlFileCount{} 个 RTL 源文件（\DesignUnitCount{} 个设计单元），目标器件为中科亿海微 EQ6HL130（CSG484\_H 封装，顶层使用 \Pins{} 个引脚），工作于单一 125~MHz 主时钟，应用侧采用 1026 答疑规定的 2 字节大端 VLID 流接口，网络侧通过 A、B 两路八位 GMII 接口收发以太网帧。验证工作涵盖 RTL 功能与性能仿真、58 份 UCDB 代码覆盖率分析、\code{tools/build_1030.py} 逻辑综合与布局布线后全量静态时序分析，以及基于 VQM 门级网表的 4 项后仿真签核（归档测试证据见附录~\ref{app:evidence}）。测试依据包括：
\begin{itemize}
  \item 竞赛选题指南（中科亿海微）及官方问题汇总 1022 / 1026 / 1030 版（涵盖 R95--R97 时延分段与同 VL 抖动定义、R110--R112 边界说明、R119 顶层 2B-VLID 应用接口、R139 1--\MaxMessageBytes{}~B IPv4 分片与去重后按序重组说明）；
  \item 协议与标准规范：ARINC 664 Part 7、RFC 791（IPv4 分片与重组）、ARINC 615A-2 / 665-1、RFC 1157 / 3416 / 3418（SNMP）、RFC 4330（SNTPv4）及 RFC 1350（TFTP）。
\end{itemize}

\section{报告结构}

第~\ref{ch:env}~章介绍测试环境、双端双网测试平台、R95/R97 与顶层 2B-VLID 流时间戳测量原理及通过判据；第~\ref{ch:strategy}~章给出 52 个需求测试点的追踪矩阵、用例清单及全设计 \DesignUnitCount{} 个设计单元的代码覆盖率分析；第~\ref{ch:func}~章分项阐述五项核心功能、顶层流包装与 1030 分片重组及联合测试的激励设计与验证结果；第~\ref{ch:perf}~章给出顶层 63-VL 声明满载规格、超规格 77/81-VL 与变长压力基线、帧格式与分片边界、分段时延与 R97 同 VL 抖动、突发缓存排空及 VL 容量的定量测量结果；第~\ref{ch:netlist}~章给出 eLinx 综合资源、布局布线后静态时序分析及 4 项门级网表仿真结果；第~\ref{ch:conclusion}~章总结测试结论与已验证边界。附录 A 至 D 依次提供完整回归明细、一键复现命令、测试证据索引及 ModelSim 原始波形截图。
